\section{Conclusion}
\label{sec:conclusion}

This paper decomposed the ``artist'' persona advantage over the ``puzzle designer'' baseline into six constituent cognitive traits, tested each trait through three phases of controlled experimentation, and identified the minimum sufficient trait set for artist-level puzzle authoring quality. Three findings constitute the core result. First, experience-first thinking (T1) is the necessary but not sufficient condition for puzzle quality: its removal causes a 29.4-point categorical failure, producing not a degraded puzzle but an absent one---a list of AoE2 terms with no clue layer, no mechanism, and no solver task. Second, the minimum sufficient set is T1 + T4 + T5 (experience-first, elegance drive, rigor/verification), which crosses the 33/45 quality threshold at Phase~3 reconstruction step R3 (35.7/45). Tier~3 traits (visual thinking, surprise instinct, systematic thinking) each contribute 4--8 additional points above the minimum sufficient baseline but are not structurally irreplaceable. Third, domain labels activate template-matching; artistic orientation activates experience-design. The puzzle designer persona produced an acrostic with an extraction table because those are the conventions the label activates. The artist persona produced a damaged-inscription mechanism because it asked first what the solver would \emph{do}, not what structure would yield the answer---and that reorientation is instantiated by the experience-first trait.

The practical application of this finding is a three-sentence creative quality injection that generalizes across domains: ``Think first about how the other person will experience this. Remove anything that is not necessary. Verify every claim before finishing.'' These sentences do not encode domain knowledge; they encode a way of working. The first reorients from content delivery to experience design, the second enforces reduction discipline, and the third provides implementation hygiene that keeps the other two honest. We predict that prepending these three sentences to any creative generation prompt---marketing copy, product design specifications, legal plain-language drafting, training material authoring, interactive fiction, game level design---will replicate the threshold-crossing effect documented here. The artist knows not what to make but how to make it for someone; these three sentences inject that knowledge operationally, without reference to any creative domain's conventions. The 15.6-point gap between domain-labeled and artistically-oriented AI output is not a puzzle-specific finding; it is a finding about the cognitive architecture of quality-oriented creative work, and the three-sentence injection is its operational form. These cross-domain predictions are grounded in the universal structure of creative receiver relationships but require empirical validation beyond puzzle hunt design before operational deployment.

Three directions merit future investigation. Cross-domain replication would test whether T1+T4+T5 produces the same threshold-crossing effect in business document generation, interactive fiction authoring, game level design, and other creative tasks where audience experience is the primary design constraint. A replication finding in two or more non-puzzle domains would establish the generalization claim on empirical rather than theoretical grounds. Multi-trait interaction analysis would test whether T1$\times$T4 produces supralinear quality gains---the elegance paradox observed here (T4 below baseline in isolation, load-bearing in combination) suggests that trait dependencies may produce non-additive effects, and identifying which trait pairs interact positively or negatively would provide a more complete causal account of the artist advantage than additive decomposition alone provides. Human validation would ask whether human puzzle designers who explicitly apply the three-trait stance to their work self-report improved output quality, and whether the C8 panel's quality judgments of human-authored puzzles align with the AI-authored findings reported here. The natural forward pointer is to the human-AI calibration question addressed in the companion paper~\cite{dellaLibera2025panel}: whether T1+T4+T5 closes the residual gap between AI-authored and human-authored puzzle quality on the subjective dimensions (Fun, Elegance, Reading Reward) where the artist persona's ceiling (42.3/45) still falls slightly short of the best human-authored conditions in that study is the question this work opens.
